This draft reports a curated Day-1 benchmark rather than a final large-scale evaluation. The paper-facing split now includes all four actions and six slices, including ambiguous-intent and insufficient-evidence cases. The strongest empirical evidence is still concentrated on answer-vs-challenge boundaries, stale-premise correction, and retrieval-conflict calibration, where source diversity and model-analysis depth are most mature. Claims about \texttt{ask} and hard \texttt{abstain} should therefore remain slice-scoped rather than presented as fully balanced evidence across all possible clarification and abstention settings. The utility metric is intended as a diagnostic summary and is always reported alongside action accuracy, over-answer rate, per-slice metrics, and confusion patterns; claims should rest on agreement between the summary score and these disaggregated behaviors.

The slice-balanced 600-example variant is a stress split, not the canonical benchmark. It adds answerable-control and conflicting-evidence rows to probe the residual slice imbalance in the 560-example split, but it also changes the answer count and is therefore used only for sensitivity analysis. Headline benchmark claims remain tied to the canonical 560-example split unless a future version re-validates and freezes a new canonical composition.

The bootstrap intervals are local uncertainty checks over the current split samples. They do not capture uncertainty from example construction, prompt-policy choice, or the decision to treat the 600-example variant as stress evidence. We therefore use them to guard against over-reading point estimates rather than as proof of a stable model ordering.

The DeepSeek reasoning runs also show low JSON adherence. We report this as part of the observed assistant behavior under the protocol, because the task requires an explicit action decision, but broader claims about reasoning models should remain scoped to this prompt and parsing setup until additional format-control runs are complete.

This draft also focuses on action choice before response generation, not on external tool-use policy. Tool-augmented behaviors can change error recovery patterns in practice, but they are not benchmarked in the current split and should be treated as future expansion work rather than implied evidence in this paper\citep{schick-etal-2023-toolformer}.

Human validation supports the active labels and paper-facing arithmetic under the documented ontology. It does not prove that the ontology is the only possible formulation of assistant action choice. Borderline examples and alternate taxonomies remain important future work.

Responsible deployment also requires care. A system that overuses \texttt{challenge} or \texttt{abstain} can become obstructive, especially for users asking ordinary answerable questions. Conversely, a system that overuses \texttt{answer} can reinforce false or stale assumptions. We therefore treat the benchmark as an evaluation of action calibration, not as a prescription that assistants should be more confrontational in all settings.

The released artifacts should avoid exposing private user data: all current examples are converted from benchmark or curated public-source material. Before submission, the anonymized package should still be checked for local paths, model-cache paths, and process notes that are useful internally but inappropriate for a double-blind review package.
